\clearpage
\chapter*{第212次报告勘误}
\markboth{第212次报告勘误}{第212次报告勘误}
\addcontentsline{toc}{chapter}{第212次报告勘误}
\sid{e212}
\setcounter{footnote}{2}
\src{301}

\noindent\textbf{5. TDTE III：商预概形。}

\begin{center}
[布尔巴基讨论班，第13卷，1960/61年度，第212号，20页。]
\end{center}

\noindent\textbf{第15页，断言 (i)。} --- 关于 $Y=X/R$ 在 $S$ 上拟射影
这一事实，目前看来只在 $R$ 来自一个等价关系时得到证明。

\noindent\textbf{第18、19页。} --- 正如我们在下一篇报告末尾所指出的，
这里所考察的猜想终究是错误的。评注8.1中指出的“肯定性事实”，似乎已由多位
作者同时证明（Nagata、Rosenlicht、Grothendieck，\ldots）。

还应指出，TDTE V 把这里发展的取商理论用于构造 Picard 概形；这一做法也可
由适当地使用 Hilbert 概形来替代（参见 SHMT）。正如第8节所说，这里所述
\src{302}
理论最重要的缺口，是缺少一个判定按非固有等价关系所取之商存在的判据，例如
由射影群的某些作用产生的等价关系。一个沿此方向的重要定理由 Mumford
得到\footnote{MUMFORD (David). --- An elementary theorem in geometric
invariant theory, \emph{Bull. Amer. Math. Soc.}, t. 67, 1961,
p. 483--486.}。关于其结果的一项改进以及对该理论的多种应用，见 SHMT。

\par
